\begin{figure}[H]
    \centering
    \begin{tikzpicture}
        \begin{axis}[
            width=0.85\textwidth,
            height=6.5cm,
            ybar,
            bar width=15pt,
            enlarge x limits=0.2,
            legend style={at={(0.05,0.95)}, anchor=north west, legend columns=1, font=\small},
            ylabel={系统内核调度构建时间 (ms)},
            xlabel={原状态测试集边数 ($|E_G|$)},
            symbolic x coords={10K, 20K, 50K, 100K},
            xtick=data,
            nodes near coords,
            every node near coord/.append style={font=\footnotesize},
            ymin=0, ymax=290,
        ]
        
        \addplot[fill=blue!40] coordinates {(10K, 15.6) (20K, 35.8) (50K, 112.4) (100K, 235.5)};
        \addplot[fill=red!40] coordinates {(10K, 0.49) (20K, 0.51) (50K, 0.94) (100K, 2.18)};
        
        \legend{CPU串行推演结构 (动态挂载), GPU前缀和无锁映射 (CSR预配)}
        \end{axis}
    \end{tikzpicture}
    \caption{线图网络转换期：动态申请模式与并发一维无锁写入体系的高频阻抗耗时量化对比}
    \label{fig:exp_phase1_topology}
\end{figure}
